\chapter{Project Management Plan}\label{ch:pm}

\section{Project Organization}

The project was carried out by a four-member team for the course teachers Mir Md.\ Saki
Kowsar and Md.\ Refaj Hossan, Department of CSE, CUET. The team was led by Turjo (Al
Rubayet). Work was divided by role so that the backend, the frontend, the database and
testing could move forward in parallel. The roles were, using the nicknames the team uses
internally:
\begin{itemize}
  \item \textbf{Turjo (Al Rubayet, 2204002)} -- team lead and frontend developer.
  \item \textbf{Saimun (Imteajul Islam, 2204015)} -- project idea, initial plan and backend
  developer.
  \item \textbf{Tahsin (Tahsin Ahmed Chowdhury, 2204028)} -- database design.
  \item \textbf{Rony (Md.~Rony Hossain, 2204005)} -- unit and integration testing,
  real-user testing and interface work.
\end{itemize}
Saimun drafted the initial plan; the whole team then sat together, finalised the plan and
agreed the future work.

\paragraph{Communication and coordination.}
All four members live in the university hall. The team met in person every week before
the lab session to plan the next steps, work together and check each other's work; no
formal minutes were recorded, and Appendix~D gives a progress log reconstructed from the
Git history. All code was shared through one GitHub repository (\texttt{Saimun-jd/RecallAI},
branch \texttt{main}). Most work was committed directly to \texttt{main}; one change (a fix
to the pnpm workspace configuration for the release build) arrived as pull request~\#1 from
a GitHub Copilot agent branch and was merged on 5~November~2026 (\Cref{ch:arch} describes
the Git workflow).

\subsection{Individual Contribution to the Project}

\Cref{tab:contrib} summarises the work distribution, and the paragraphs below describe each
member's contribution. Of the 48 commits on \texttt{main} between 22~July and
5~November~2026, 31 were made from the GitHub account \texttt{Saimun-jd} (Saimun), 14 from
\texttt{rubayet2027} and one from \texttt{Al Rubayet} (both Turjo's account; the latter is
the pull-request merge), and two came from the Copilot agent. Most of the
\texttt{rubayet2027} commits, including the backend work of September, were made on
Saimun's computer while it was logged in to Turjo's account, so the commit author does not
show who wrote the code. Commit counts also do not measure planning, design, testing or
documentation work.

\begin{table}[htbp]
\centering
\caption{Individual contribution to the project}\label{tab:contrib}
\scriptsize\setlength{\tabcolsep}{2.2pt}
\begin{tabular}{|l|c|c|c|c|c|c|c|c|}
\hline
\textbf{Member} & \makecell{\textbf{Req.}\\\textbf{Spec.}} & \textbf{Planning} & \textbf{Design} &
\makecell{\textbf{AI / Model}\\\textbf{Integration}} & \makecell{\textbf{User}\\\textbf{Interface}} & \textbf{Testing} &
\textbf{Deployment} & \textbf{Docs}\\
\hline
Al Rubayet (Turjo)               & \cmark & \cmark & \cmark &        & \cmark &        &        & \cmark \\ \hline
Imteajul Islam (Saimun)          & \cmark & \cmark & \cmark & \cmark &        &        & \cmark & \cmark \\ \hline
Tahsin Ahmed Chowdhury (Tahsin)  & \cmark & \cmark & \cmark &        &        & \cmark &        & \cmark \\ \hline
Md.~Rony Hossain (Rony)          & \cmark & \cmark &        &        & \cmark & \cmark &        & \cmark \\ \hline
\end{tabular}
\end{table}

\begin{itemize}
  \item \textbf{Al Rubayet (Turjo, 2204002).} Led the team and developed the frontend of
  the desktop application. Also took part in the requirements, planning, design and
  documentation work.
  \item \textbf{Imteajul Islam (Saimun, 2204015).} Proposed the project idea, drafted the
  initial plan and wrote the backend code, including the AI provider integration. Also
  handled deployment (the Windows release) and took part in the requirements, design and
  documentation work.
  \item \textbf{Tahsin Ahmed Chowdhury (Tahsin, 2204028).} Designed the database and took
  part in the requirements, planning, design, testing and documentation work.
  \item \textbf{Md.~Rony Hossain (Rony, 2204005).} Carried out the unit and integration
  testing, organised the real-user testing sessions and worked on the user interface. Also
  took part in the requirements, planning and documentation work.
\end{itemize}

\section{Process Model Used}

The team used an \textbf{incremental model with iterative refinement}
\cite{pressman2019,sommerville2016}. A first, narrow version of the core idea (turning a PDF
into on-demand flashcards) was built as a command-line prototype. The system was then
grown in increments. Each increment went through design, implementation, testing against
documents the team had uploaded, and integration into \texttt{main}, and each produced a
version of the product that could be run and demonstrated. Later increments also revisited
earlier ones: the interface was redesigned twice in August, and the study features built
in August were reconciled with the new \texttt{/api/v1} platform in September.

The increments, dated from the repository history, are listed in \Cref{tab:increments},
and \Cref{fig:process} shows how the model was applied.

\begin{table}[htbp]
\centering
\caption{Increments of RecaLL, dated from the Git history}\label{tab:increments}
\small
\begin{tabularx}{\textwidth}{@{}l >{\raggedright\arraybackslash}p{2.3cm} X r@{}}
\toprule
\textbf{Inc.} & \textbf{Dates (2026)} & \textbf{Main deliverables} & \textbf{Commits}\\
\midrule
1 & 22--24 Jul & Python command-line prototype: PDF text, code and diagram extraction with
source pages; on-demand flashcards with topic similarity; hardened API and storage & 9\\
2 & 5 Aug & Tauri desktop application and Notion-style Markdown note editor; hard-coded
API keys removed & 3\\
3 & 19--25 Aug & Socratic drill learning mechanism; LaTeX rendering fixes; two interface
redesigns; release workflow and first Windows executable (23~Aug); pre-release
\texttt{v0.1.5} (24~Aug); sidecar packaging fix & 10\\
4 & 26 Aug--9 Sep & Supabase integration and sign-in flow; Marker OCR extraction; chat with
an AI teacher and topic-specific chat; per-concept drills; new notes feature; chunked
processing of long documents & 7\\
5 & 12--21 Sep & \texttt{/api/v1} core: document ingestion pipeline, grounded chat,
flashcard sets, quizzes and attempts, spaced review; reconcile merge with the earlier study
features (16~Sep); review-queue interface; backend hardening & 11\\
6 & 25--29 Sep & Production hardening, tenant-isolation test matrix, container files,
\texttt{ci.yml} (25~Sep), code clean-up, Knowledge Hub and interface fixes & 4\\
7 & October & University closed; no project work and no commits & 0\\
8 & 5 Nov & Review and stabilisation: fixes for the security, user-experience and logic
findings of the code review; PR~\#1 merged; optional Supabase sync made skippable
& 4\\
\bottomrule
\end{tabularx}
\end{table}

\figplaceholder[0.9\textwidth]{process_model.pdf}{Incremental process model as applied to
RecaLL}{fig:process}

\subsection{Rationale for Choosing Lifecycle Model}

\begin{itemize}
  \item \textbf{Project size and modularity.} The product splits into loosely coupled
  features (import, reader, notes, drill, flashcards, review, chat, quizzes, settings).
  Each can be built and tested on its own and added to a running product, which is what
  increments are designed for.
  \item \textbf{Requirement stability.} The core idea (PDF to active recall) was stable from
  the start, but the details were not: on-demand card generation replaced batch
  generation on 24~July, the drill was redesigned per concept in September, and the
  interface changed several times. Increments allowed the stable core to be built early and
  the volatile parts to change later without re-planning the whole project.
  \item \textbf{Team size and skills.} A four-person student team, sharing the term with
  other courses, cannot run the full ceremony of Scrum (product owner, scrum master, sprint
  reviews and retrospectives), but still benefits from frequent working builds.
  \item \textbf{Risk.} The riskiest technical parts (PDF extraction, LLM output quality and
  desktop packaging) were tried first, in the prototype and the first three increments. An
  installer existed by 23~August, so packaging problems were found early rather than at the
  end of the term.
  \item \textbf{Customer involvement.} The course teachers act as the customer and see progress
  in the sessional classes. A working increment is easy to demonstrate.
\end{itemize}

\paragraph{Rejected alternatives.} \emph{Waterfall} was rejected because it produces no
working software until late in the term and copes badly with the interface and prompt
changes that this project needed. \emph{Spiral} requires a formal risk analysis in every
loop, which is too heavy for a student project. \emph{Pure Scrum} was rejected for the
ceremony overhead described above. \emph{Throw-away prototyping} would have discarded the
command-line prototype, whose extraction and flashcard code was carried into the desktop
application instead.

\section{Risk Analysis}

Risk exposure is computed as probability $\times$ impact on a 1--3 scale (Low~=~1,
Medium~=~2, High~=~3). \Cref{tab:risk} lists the risks identified, their exposure and the
mitigation used in the delivered system.

\begin{small}\setstretch{1.1}
\begin{longtable}{@{}l >{\raggedright\arraybackslash}p{3.6cm} c c c >{\raggedright\arraybackslash}p{6.0cm}@{}}
\caption{Risk analysis}\label{tab:risk}\\
\toprule
\textbf{ID} & \textbf{Risk} & \textbf{Prob.} & \textbf{Impact} & \textbf{Exp.} & \textbf{Mitigation}\\
\midrule
\endfirsthead
\toprule
\textbf{ID} & \textbf{Risk} & \textbf{Prob.} & \textbf{Impact} & \textbf{Exp.} & \textbf{Mitigation}\\
\midrule
\endhead
\bottomrule
\endfoot
R1 & AI provider unavailable, slow or out of quota & High & Med & 6 & Four providers
including local Ollama; per-provider rate limiter (for example 5 requests/min on the Gemini
free tier); retry with exponential back-off on HTTP 429/502/503/504\\
R2 & AI produces wrong or invented content in notes, cards or answers & Med & High & 6 &
Prompts ground the model in the source text; citations checked against retrieved chunks;
cards and notes link to their source page and remain editable\\
R3 & LLM returns malformed JSON & High & Med & 6 & \texttt{json-repair}, a truncation-repair
pass and Pydantic schema validation before anything is saved\\
R4 & Prompt injection through document content & Med & High & 6 & Document text is wrapped in
\texttt{<reference\_data>} blocks, closing tags are sanitised, and the system prompt tells
the model never to follow instructions inside them\\
R5 & Leakage of personal API keys & Med & High & 6 & AES-256-GCM vault with per-install key;
only the last four characters are shown (\Cref{sec:security})\\
R6 & Another program on the computer calls the local backend & Med & High & 6 &
Backend binds to \texttt{127.0.0.1} only; CORS allow-list of the local origins\\
R7 & Two databases (application and per-user study) drift apart & Med & Med & 4 &
\texttt{BookSyncService} mirrors books into workspaces; upload writes both stores in one
request\\
R8 & Scanned or handwritten PDFs yield no text & Med & Med & 4 & Optional Marker extractor
(local or through the Datalab API) on the book and topic path\\
R9 & Desktop packaging of the Python sidecar fails & Med & High & 6 & Installer built from the
third increment on; release workflow on \texttt{windows-latest} builds the PyInstaller
sidecar and runs \texttt{tauri-action} on every \texttt{v*} tag\\
R10 & Team member unavailable (examinations, illness) & Med & High & 6 & Shared repository,
README and diagram documents so another member can continue a module\\
R11 & Schedule slip near the end of term & High & Med & 6 & Incremental delivery keeps a
working build; scope cut by priority (real payments and cloud sync left out)\\
R12 & Dependency drift breaks the build (\texttt{requirements.txt} is unpinned except
\texttt{json-repair}) & Med & Med & 4 & CI installs and tests on every push; lock file for
the frontend (\texttt{pnpm-lock.yaml}). Planned: pin backend versions\\
\end{longtable}
\end{small}

\section{Constraints to Project Implementation}

\begin{itemize}
  \item \textbf{Schedule.} The project ran from 20~July~2026 (the first commit was on
  22~July) to the submission on 19~November~2026, shared with other courses and
  examinations. The university was closed in October 2026, so no project work was done
  that month.
  \item \textbf{Budget.} No funding. Only free and open-source tools, free API tiers and the
  team's own laptops could be used (\Cref{sec:budget}).
  \item \textbf{Hardware.} Development on personal laptops without a dedicated GPU, which
  rules out training a model and limits the size of a local Ollama model.
  \item \textbf{AI services.} Cloud providers impose request quotas (the free Gemini tier is
  configured at 5 requests per minute) and need an internet connection; output quality and
  latency vary between providers and models.
  \item \textbf{Data.} By the project's own rule (\texttt{AGENTS.md}), the system is tested
  only with documents the users have actually uploaded, never with synthetic documents. This
  keeps the tests realistic but limits the size of the test set.
  \item \textbf{Platform.} The release workflow builds a Windows installer, the platform the
  team develops and tests on. Tauri can also target macOS and Linux.
  \item \textbf{Legal.} Students' documents may be copyrighted; RecaLL stores them only on the
  student's computer. Personal data must be handled according to Bangladesh's data-protection
  rules, the Personal Data Protection Ordinance 2025 (PDPO 2025) \cite{pdpo2025}, and, for any EU user, the
  GDPR \cite{gdpr2016}. Each AI provider's terms of use apply to the text sent to it.
  \item \textbf{Skills.} The team had to learn FastAPI, Tauri~v2, React~19, prompt design and
  FSRS scheduling during the project.
\end{itemize}

\section{Tools/Languages Requirements}

\subsection{Hardware Requirements}

\Cref{tab:hw} lists the hardware used for development and build, and the hardware a student
needs to run RecaLL.

\begin{table}[htbp]
\centering
\caption{Hardware requirements}\label{tab:hw}
\small
\begin{tabularx}{\textwidth}{@{}>{\raggedright\arraybackslash}p{3.0cm} X X@{}}
\toprule
\textbf{Purpose} & \textbf{Minimum} & \textbf{Used / recommended}\\
\midrule
Development machines & 64-bit CPU, 8~GB RAM, 10~GB free disk & Four Windows laptops with a
512~GB SSD each: Intel Core~i5 with 32~GB RAM (Saimun); AMD Ryzen~5 with 16~GB (Turjo); AMD
Ryzen~7 with 16~GB (Rony); Intel Core~i5 with 16~GB (Tahsin)\\
CI build & GitHub-hosted runner & \texttt{ubuntu-latest} (tests, lint, build, E2E)\\
Release build & GitHub-hosted Windows runner & \texttt{windows-latest} (PyInstaller sidecar,
Tauri installer)\\
Student computer & 64-bit Windows; the backend used 119~MB of RAM idle and at most 381~MB
during AI processing; the team's library of 11 documents occupies about 120~MB of disk &
16~GB RAM and an SSD, as on the team's laptops\\
Local AI (optional) & Additional RAM for the chosen Ollama model; default
\texttt{gemma3:4b} & A GPU speeds up generation but is not required\\
Network & Needed only for cloud AI providers and optional sync & Offline use with Ollama\\
\bottomrule
\end{tabularx}
\end{table}

\subsection{Software Requirements}

\Cref{tab:sw} lists the software stack. Backend versions are the ones installed in the
project's virtual environment (\texttt{requirements.txt} does not pin most versions), and
frontend versions are the ones installed from \texttt{desktop/package.json} and its lock
file.

\begin{small}
\begin{longtable}{@{}>{\raggedright\arraybackslash}p{2.8cm} >{\raggedright\arraybackslash}p{5.8cm} >{\raggedright\arraybackslash}p{5.4cm}@{}}
\caption{Software requirements (installed versions)}\label{tab:sw}\\
\toprule
\textbf{Category} & \textbf{Tool / library (version)} & \textbf{Use}\\
\midrule
\endfirsthead
\toprule
\textbf{Category} & \textbf{Tool / library (version)} & \textbf{Use}\\
\midrule
\endhead
\bottomrule
\endfoot
Languages & Python 3.13.5 (local), 3.12 (CI); TypeScript 6.0.3; Rust (edition 2021); SQL &
Backend, frontend, desktop shell, schema\\
Backend framework & FastAPI 0.162.0, Starlette 1.10.0, Uvicorn 0.54.0 & REST API and ASGI
server\\
Validation / config & Pydantic 2.14.2, pydantic-settings 2.16.2 & Request schemas,
environment settings\\
Documents & PyMuPDF 2.1, pymupdf4llm 2.1; Marker (optional) & PDF parsing, Markdown
extraction, OCR\\
Scheduling & fsrs 6.3.2 & FSRS-6 spaced repetition\\
AI and HTTP & httpx 0.28.1, json-repair 0.63.3, numpy 2.5.4, langfuse 5.1.0 (optional) &
Provider calls, JSON repair, vector maths, tracing\\
Security & PyJWT 2.15.1, cryptography 50.0.2 & HS256 tokens, AES-256-GCM\\
Storage and sync & SQLite (WAL mode), platformdirs 4.14.3, supabase 2.33.0 (optional) &
Local databases, data folder, cloud sync\\
Frontend & React 19.2.8, react-router-dom 7.18.1, Redux Toolkit 2.12.0, Tailwind CSS 4.3.3,
Vite 8.1.5 & Interface, routing, state, styling, build\\
Frontend libraries & KaTeX 0.18.1, pdfjs-dist 6.1.200, Recharts 3.10.1, BlockNote 0.52.1 &
Maths, PDF viewer, charts, note editor\\
Desktop & Tauri 2.11.3 (core), @tauri-apps/cli 2.11.4, @tauri-apps/api 2.11.1 & Desktop
shell, sidecar, installer\\
Packaging & PyInstaller 6.22.3 & Backend sidecar executable\\
Testing and quality & pytest 9.2.0, Playwright 1.62.1, oxlint 1.75.0 & Unit and integration
tests, E2E journeys, linting\\
Build tools & Node.js 22.22.2 (local), 20 (CI); pnpm 12.18.0 (local), 9 (CI) & Frontend
toolchain\\
AI providers & Gemini, OpenAI, Groq, OpenRouter (cloud); Ollama (local) & Language models\\
DevOps and collaboration & Git 2.47.1, GitHub, GitHub Actions, GitHub Copilot agent & Version
control, CI/CD, one automated pull request\\
Editors and IDEs & Antigravity (with its AI agent) and VS Code & Development\\
\end{longtable}
\end{small}

\section{Project Timeline and Schedule}

\Cref{fig:gantt} shows the project timeline in calendar weeks from 20~July to
22~November~2026. The bars for the increments are taken from the dates of the commits in
\Cref{tab:increments}. Testing was done alongside each increment rather than as a separate
phase; the final automated test run used in \Cref{ch:test} was on 16~November~2026. The
university was closed in October, and the report was submitted on 19~November~2026.

\begin{figure}[htbp]
\centering
\resizebox{\textwidth}{!}{%
\begin{ganttchart}[hgrid, vgrid={*{1}{draw=black!30, dotted}}, x unit=0.62cm,
  y unit title=0.6cm, y unit chart=0.62cm,
  title label font=\scriptsize, bar label font=\small, milestone label font=\small\itshape,
  title/.append style={fill=white, draw=black},
  bar/.append style={fill=black!15, draw=black},
  milestone/.append style={fill=black, draw=black}]{1}{18}
\gantttitle{2026 (week starting Monday)}{18}\\
\gantttitle{Jul}{2}\gantttitle{Aug}{5}\gantttitle{Sep}{4}\gantttitle{Oct}{4}\gantttitle{Nov}{3}\\
\gantttitlelist{1,...,18}{1}\\
\ganttbar{Idea and planning}{1}{1}\\
\ganttbar{Inc.~1: CLI prototype}{1}{1}\\
\ganttbar{Inc.~2: desktop app, notes}{3}{3}\\
\ganttbar{Inc.~3: drill, new UI, release}{5}{6}\\
\ganttbar{Inc.~4: login, OCR, AI chat}{6}{8}\\
\ganttbar{Inc.~5: \texttt{/api/v1} core}{8}{10}\\
\ganttbar{Inc.~6: hardening, CI}{10}{11}\\
\ganttbar{University closed (no work)}{12}{15}\\
\ganttbar{Inc.~8: review fixes}{16}{16}\\
\ganttbar{Final test run and report}{18}{18}\\
\ganttmilestone{Pre-release v0.1.5 (24 Aug)}{6}\\
\ganttmilestone{\texttt{ci.yml} added (25 Sep)}{10}\\
\ganttmilestone{PR \#1 merged (HEAD 60d40f9)}{16}\\
\ganttmilestone{Submission (19 Nov)}{18}
\end{ganttchart}}
\caption{Project timeline (Gantt chart). Week 1 starts on 20~July~2026}\label{fig:gantt}
\end{figure}

\section{Estimated Budget}\label{sec:budget}

RecaLL is a desktop application that runs on the student's own computer, so it needs no
server, domain or hosting. No money was spent on the project: every tool and service the
team used was free. Some optional AI providers in the application need a paid key, but the
team did not use them. \Cref{tab:budget} lists the items.

\begin{table}[htbp]
\centering
\caption{Project budget}\label{tab:budget}
\small
\begin{tabularx}{\textwidth}{@{}X l r r@{}}
\toprule
\textbf{Item} & \textbf{Unit} & \textbf{Qty} & \textbf{Cost (BDT)}\\
\midrule
Development tools (Python, Node.js, Tauri, Antigravity, VS Code) & -- & -- & 0\\
GitHub repository and GitHub Actions minutes (free tier) & month & 4 & 0\\
Hosting, domain and TLS certificate (not needed for a desktop app) & -- & -- & 0\\
Local AI through Ollama & -- & -- & 0\\
Gemini API, free tier & month & 4 & 0\\
Paid AI provider keys (optional in the app; not used by the team) & -- & -- & 0\\
Supabase (optional) & -- & -- & 0\\
Code-signing certificate for the installer (not purchased) & year & 0 & 0\\
\midrule
\multicolumn{3}{@{}l}{\textbf{Total cash cost}} & \textbf{0}\\
\bottomrule
\end{tabularx}
\end{table}

For a student using RecaLL, the running cost is zero with Ollama or a free API tier. With a
paid provider, the student pays that provider directly through their own key (BYOK); RecaLL
charges nothing and counts such requests as zero credits.

\section{Technological, Societal, \& Environmental Impact of the Project}

\paragraph{Positive impacts.}
\begin{itemize}
  \item \textbf{Individuals.} Students spend less time on the mechanical parts of studying
  (structuring, note drafting, card writing, scheduling) and more on retrieval practice and
  spaced revision, the two techniques with the strongest evidence \cite{dunlosky2013}. The
  Socratic drill shows a student \emph{what} they misunderstood, not only that an answer was
  wrong.
  \item \textbf{Privacy and ownership.} Study material, notes and progress stay on the
  student's computer. With Ollama, no document text leaves the machine at all.
  \item \textbf{Cost and access.} The product is free to run with a local model or a free API
  tier, which matters for students who cannot pay for subscription study tools.
  \item \textbf{Technology.} The project shows a reusable pattern for student teams: a local
  Python AI backend packaged as a sidecar inside a small Tauri desktop application, with a
  provider-neutral AI layer and a CI-built installer.
\end{itemize}

\paragraph{Negative impacts and mitigations.}
\begin{itemize}
  \item \textbf{Over-reliance and academic integrity.} Students may accept AI-written notes
  without checking them, or use generated material in assessed work. Notes are editable,
  cards and answers link to their source page, and the drill requires the student to write
  their own explanation.
  \item \textbf{Errors in AI output.} A model can be confidently wrong, and a wrong flashcard
  is then revised repeatedly. Citations are checked against retrieved chunks, and every card
  can be edited or deleted. \Cref{ch:sustain} discusses this further.
  \item \textbf{Energy use.} Every cloud LLM call uses data-centre energy, and a local model
  uses the student's own power. RecaLL limits context (five chunks, at most 3{,}000 tokens of
  reference text and six history messages per chat turn) and generates content on demand
  rather than for a whole book at once.
  \item \textbf{Digital divide.} Cloud models need an internet connection; the application
  itself is free, and Ollama allows offline use.
  \item \textbf{Copyright.} Uploaded textbooks are processed locally and never shared; their text
  reaches a cloud provider only when an AI feature uses one.
\end{itemize}
